\section{可数性公理}\label{section:ch04-countability-axioms}

首先，我们重新陈述第 21 节中给出过的一个定义。

\begin{definition}\label{definition:ch04-first-countable}
空间 $X$ 称为在 $x$ 处有\textbf{可数基}（countable basis at $x$），如果存在 $x$ 的邻域的一个可数族 $\mathcal{B}$，使得 $x$ 的每一个邻域都至少包含 $\mathcal{B}$ 中的一个成员。如果在空间的每一点处都有可数基，则称这个空间满足\textbf{第一可数性公理}（first countability axiom）或\textbf{第一可数的}（first-countable）。
\end{definition}

我们已经注意到，每一个可度量化空间都满足这个公理，见第 21 节。

考虑满足这个公理的空间的一个重要原因就是，这类空间有一个非常有用的特性，即只要借助于这类空间中的收敛序列就可以确定集合的极限点，并且可以验证函数的连续性。以前我们已经注意到了这一点，现在正式地作为定理来叙述它：

\begin{theorem}\label{theorem:ch04-sequence-characterization}
设 $X$ 是一个拓扑空间。
\begin{enumerate}
\item 设 $A$ 是 $X$ 的一个子集。若存在 $A$ 中点的序列收敛到 $x$，则 $x\in\overline{A}$。若 $X$ 满足第一可数性公理，则其逆命题也成立.
\item 设 $f\colon X\to Y$。若 $f$ 是连续的，则对 $X$ 中的每一个收敛序列 $x_{n}\to x$，序列 $f(x_{n})$ 收敛于 $f(x)$。如果 $X$ 满足第一可数性公理，那么其逆命题也成立.
\end{enumerate}
\end{theorem}

这个定理的证明实际上是第 21 节中在可度量化假设下所给证明的一个直接推广，在此不再重新陈述。

下述公理比第一可数性公理更为重要：

\begin{definition}\label{definition:ch04-second-countable}
若拓扑空间 $X$ 具有可数基，则称 $X$ 满足\textbf{第二可数性公理}（second countability axiom）或称 $X$ 为\textbf{第二可数的}（second-countable）。
\end{definition}

显然第二可数性公理蕴涵着第一可数性公理：若 $\mathcal{B}$ 是 $X$ 的一个可数拓扑基，则 $\mathcal{B}$ 中由包含着点 $x$ 的那些元素所组成的子族就是 $x$ 处的一个可数基。事实上，第二可数性公理比第一可数性公理强得多，甚至并不是每一个度量空间都能满足它。

那么，为什么要对这个性质感兴趣呢？因为，一方面，许多熟知的空间具有这个性质。另一方面，正如我们将要看到的，它是证明像 Urysohn 度量化定理这样一些定理时要用到的一个十分重要的假设。

\begin{example}\label{example:ch04-second-countable-examples}
实直线 $\mathbb{R}$ 具有可数基，即所有端点为有理数的开区间 $(a,b)$ 的族。同样，$\mathbb{R}^{n}$ 具有可数基，即所有端点为有理数的开区间的积的族。甚至 $\mathbb{R}^{\omega}$ 也有可数基，即所有积 $\prod_{n\in\mathbb{Z}_{+}}U_{n}$ 构成的族，其中对于有限多个 $n$ 而言，$U_{n}$ 是端点为有理数的开区间，而对于其他所有的 $n$，$U_{n}=\mathbb{R}$。
\end{example}

\begin{example}\label{example:ch04-uniform-not-second-countable}
在一致拓扑下，$\mathbb{R}^{\omega}$ 满足第一可数性公理（因为它是可度量化的），却不满足第二可数性公理。为了证明这一点，我们先要证明：若 $X$ 有可数基 $\mathcal{B}$，则 $X$ 的任意离散子空间 $A$ 必定是可数的。对于每一个 $a\in A$，选取基中的一个元素 $B_{a}$，使得它与 $A$ 只交于点 $a$。若 $a$ 和 $b$ 是 $A$ 中不同的两个点，则 $B_{a}$ 和 $B_{b}$ 也不同。这是因为第一个集合包含点 $a$，但第二个却不包含。由此可见，映射 $a\mapsto B_{a}$ 是从 $A$ 到 $\mathcal{B}$ 的一个单射，所以 $A$ 必是可数的。

我们注意，$\mathbb{R}^{\omega}$ 中所有 $0$ 和 $1$ 序列构成的子空间 $A$ 是不可数的，并且具有离散拓扑。因为对 $A$ 中任意两个不同的点 $a$ 和 $b$，$\overline{\rho}(a,b)=1$。因此在一致拓扑下，$\mathbb{R}^{\omega}$ 没有可数基。
\end{example}

这两个可数性公理都具有在取子空间运算或可数积运算下保持不变这一很好的性质。

\begin{theorem}\label{theorem:ch04-countability-hereditary}
第一可数空间的子空间是第一可数的。第一可数空间的可数积是第一可数的。第二可数空间的子空间是第二可数的。第二可数空间的可数积是第二可数的。
\end{theorem}

\begin{proof}
考虑第二可数性公理。若 $\mathcal{B}$ 是 $X$ 的一个可数基，则 $\{\,B\cap A\mid B\in\mathcal{B}\,\}$ 便是 $X$ 的子空间 $A$ 的一个可数基。若 $\mathcal{B}_{i}$ 是空间 $X_{i}$ 的一个可数基，则所有积 $\prod U_{i}$ 构成的族便是 $\prod X_{i}$ 的一个可数基，其中 $U_{i}$ 满足条件：对于有限多个 $i$，$U_{i}\in\mathcal{B}_{i}$。而对于其他的所有 $i$，$U_{i}=X_{i}$。

关于第一可数性公理的证明是类似的。
\end{proof}

下面定理中给出关于第二可数性公理的两个结论，它们在后面会用到。首先给出以下定义：

\begin{definition}\label{definition:ch04-dense}
空间 $X$ 的子集 $A$ 称为在 $X$ 中是\textbf{稠密的}（dense），如果 $\overline{A}=X$。
\end{definition}

\begin{theorem}\label{theorem:ch04-second-countable-lindelof-separable}
设 $X$ 有一个可数基。则
\begin{enumerate}
\item $X$ 的每一个开覆盖有一个可数字族覆盖 $X$.
\item $X$ 存在一个可数子集在 $X$ 中稠密.
\end{enumerate}
\end{theorem}

\begin{proof}
设 $\{B_{n}\}$ 是 $X$ 的一个可数基。

(a) 令 $\mathcal{A}$ 为 $X$ 的一个开覆盖。对于每一个正整数 $n$，只要可能，我们就选取一个 $A_{n}\in\mathcal{A}$，使得 $A_{n}$ 包含基元素 $B_{n}$。那么，这些集合 $A_{n}$ 构成的族 $\mathcal{A}'$ 是可数的，因为它的下标集 $J$ 是正整数的一个子集。而且 $\mathcal{A}'$ 覆盖 $X$，因为任意给定 $X$ 的一个点 $x$，可以在 $\mathcal{A}$ 中选择一个集合 $A$ 包含点 $x$。因为 $A$ 是开的，所以存在基元素 $B_{n}$ 使得 $x\in B_{n}\subset A$。因而 $B_{n}$ 就必定包含在 $\mathcal{A}$ 的一个成员中，其下标属于集合 $J$。于是 $A_{n}$ 有定义。并且 $A_{n}$ 包含 $B_{n}$，所以 $A_{n}$ 包含点 $x$。

(b) 从基中的每一个非空元素 $B_{n}$ 中选取一点 $x_{n}$。设这些点 $x_{n}$ 构成集合 $D$，那么 $D$ 在 $X$ 中稠密。这是因为对于 $X$ 中任意给定的一点 $x$，每一个包含 $x$ 的基元素都和 $D$ 相交，所以 $x$ 属于 $\overline{D}$。
\end{proof}

\cref{theorem:ch04-second-countable-lindelof-separable} 所列出的两个性质有时也分别被当作一种可数性公理。每一个开覆盖都包含可数子覆盖的空间，通常称为\textbf{Lindelöf 空间}（Lindelöf space）。有可数稠密子集的空间常被称为\textbf{可分的}（separable）。一般说来，这两个性质都比第二可数性公理弱。但若空间是可度量化的，则它们与第二可数性公理等价（见习题 5）。就其重要性而言，它们不及第二可数性公理，但也不可忽视，使用它们会给我们带来方便。例如，通常证明一个空间 $X$ 具有可数稠密子集就比证明 $X$ 具有可数基容易。如果空间还是可度量化的（像分析中常见的那样），那么这就蕴涵了 $X$ 是第二可数的。

我们并不应用这些性质去证明任何定理，然而其中的 Lindelöf 条件会在处理一些例子时用到。正如下面的一些例子所表明的，对于取子空间和笛卡儿积的运算而言，这两个性质都不像我们所希望的那样能够得以保持。

\begin{example}\label{example:ch04-sorgenfrey-line-countability}
空间 $\mathbb{R}_{\ell}$ 除了不满足第二可数性公理外，满足其他所有的可数性公理。

给定 $x\in\mathbb{R}_{\ell}$，那么所有形如 $[x,x+1/n)$ 的基元素构成的集合是 $x$ 处的一个可数基。容易看出有理数集在 $\mathbb{R}_{\ell}$ 中稠密。

为了证明 $\mathbb{R}_{\ell}$ 没有可数基，设 $\mathcal{B}$ 是 $\mathbb{R}_{\ell}$ 的一个基。对于任意 $x$，选取 $\mathcal{B}$ 中的一个元素 $B_{x}$，使得 $x\in B_{x}\subset[x,x+1)$。若 $x\neq y$，则 $B_{x}\neq B_{y}$。这是因为 $x=\inf B_{x}$，$y=\inf B_{y}$。因此 $\mathcal{B}$ 必定是不可数的。

要证明 $\mathbb{R}_{\ell}$ 是 Lindelöf 空间，还需要做一些工作。但只要证明 $\mathbb{R}_{\ell}$ 中任何一个由基元素构成的开覆盖包含着可数字族覆盖 $\mathbb{R}_{\ell}$ 就行了（此处请自行验证）。为此设
\begin{align*}
\mathcal{A}=\{\,[a_{\alpha},b_{\alpha})\,\}_{\alpha\in J}
\end{align*}
是由下限拓扑的基元素构成的 $\mathbb{R}$ 的一个覆盖。我们希望找到它的一个可数字族覆盖 $\mathbb{R}$。

令 $C$ 为一个集合，
\begin{align*}
C=\bigcup_{\alpha\in J}(a_{\alpha},b_{\alpha}),
\end{align*}
它是 $\mathbb{R}$ 的一个子集。下面证明 $\mathbb{R}\setminus C$ 是可数的。

设 $x$ 是 $\mathbb{R}\setminus C$ 的一个点。我们知道 $x$ 不属于任何一个开区间 $(a_{\alpha},b_{\alpha})$。因此，存在指标 $\beta$ 使得 $x=a_{\beta}$。选取这样的一个 $\beta$ 和区间 $(a_{\beta},b_{\beta})$ 中的一个有理数 $q_{x}$。因 $(a_{\beta},b_{\beta})$ 包含于 $C$ 中，所以区间 $(a_{\beta},q_{x})=(x,q_{x})$ 也包含于 $C$ 中。由此推出，如果 $x$ 和 $y$ 是 $\mathbb{R}\setminus C$ 中不同的两个点，并且 $x<y$，那么必有 $q_{x}<q_{y}$。（因为如若不然，我们将会有 $x<y<q_{y}<q_{x}$，于是 $y$ 必然要落在区间 $(x,q_{x})$ 中，因此也就在 $C$ 中。）从而，从 $\mathbb{R}\setminus C$ 到 $\mathbb{Q}$ 的映射 $x\mapsto q_{x}$ 是一个单射，因此 $\mathbb{R}\setminus C$ 是可数的。

现在我们证明 $\mathcal{A}$ 的某一个可数字族覆盖 $\mathbb{R}$。先对于 $\mathbb{R}\setminus C$ 中的每一个元素选取 $\mathcal{A}$ 中的一个成员包含它，这样便得到了 $\mathcal{A}$ 的一个可数字族 $\mathcal{A}'$ 覆盖 $\mathbb{R}\setminus C$。取集合 $C$，并且赋予这个集合 $\mathbb{R}$ 的子空间的拓扑。相对于这个拓扑，$C$ 满足第二可数性公理。$C$ 被这些集合 $(a_{\alpha},b_{\alpha})$ 所覆盖，它们在 $\mathbb{R}$ 中都是开集，因此在 $C$ 中也都是开集。于是存在可数字族覆盖 $C$。设这个可数字族由 $(a_{\alpha},b_{\alpha})$ 这些元素构成，其中 $\alpha=\alpha_{1},\alpha_{2},\cdots$。于是族
\begin{align*}
\mathcal{A}''=\{\,[a_{\alpha},b_{\alpha})\mid\alpha=\alpha_{1},\alpha_{2},\cdots\,\}
\end{align*}
是 $\mathcal{A}$ 的可数字族，并且覆盖 $C$，并且 $\mathcal{A}'\cup\mathcal{A}''$ 便是 $\mathcal{A}$ 的一个可数字族，覆盖 $\mathbb{R}_{\ell}$。
\end{example}

\begin{example}\label{example:ch04-sorgenfrey-plane-not-lindelof}
两个 Lindelöf 空间的积未必是 Lindelöf 的。尽管空间 $\mathbb{R}_{\ell}$ 是 Lindelöf 的，我们将证明 $\mathbb{R}_{\ell}\times\mathbb{R}_{\ell}$ 却不是 Lindelöf 的。$\mathbb{R}_{\ell}^{2}$ 空间是拓扑学中一个非常有用的例子，称之为\textbf{Sorgenfrey 平面}（Sorgenfrey plane）。

空间 $\mathbb{R}_{\ell}^{2}$ 中所有形如 $[a,b)\times[c,d)$ 的集合构成了它的一个基。为了证明它不是一个 Lindelöf 空间，我们考虑它的一个子空间
\begin{align*}
L=\{\,x\times(-x)\mid x\in\mathbb{R}_{\ell}\,\}.
\end{align*}
容易验证 $L$ 是 $\mathbb{R}_{\ell}^{2}$ 中的一个闭集。用开集 $\mathbb{R}_{\ell}^{2}\setminus L$ 和所有形如
\begin{align*}
[a,b)\times[-a,d)
\end{align*}
的基元素构成 $\mathbb{R}_{\ell}^{2}$ 的一个开覆盖。

这些开集中的每一个最多与 $L$ 交于一点。由于 $L$ 是不可数的，所以不可能有可数字族覆盖 $\mathbb{R}_{\ell}^{2}$，如\cref{figure:ch04-sorgenfrey-plane}所示。
\end{example}

\par
\begin{figure}[htbp]
\centering
\includegraphics[width=0.5\textwidth]{ch04-fig30-1.png}
\caption{Sorgenfrey 平面不是 Lindelöf 的}\label{figure:ch04-sorgenfrey-plane}
\end{figure}

\begin{example}\label{example:ch04-lindelof-subspace}
Lindelöf 空间的子空间未必是 Lindelöf 空间。有序矩形 $I_{o}^{2}$ 是紧致的，所以是 Lindelöf 的。但子空间 $A=I\times(0,1)$ 却不是 Lindelöf 的。因为 $A$ 是两两无交的集合 $U_{x}=\{x\}\times(0,1)$ 的并，其中的每一个集合都是 $A$ 中的开集。而这些集合构成的集族是不可数的，并且没有真子集族能覆盖 $A$。
\end{example}

\subsection*{习题}

\begin{enumerate}
\item (a) 若集合 $A$ 是空间 $X$ 中可数个开集的交，则称 $A$ 是 $X$ 中的一个 $G_{\delta}$ 集。证明在满足第一可数性公理的 $T_{1}$ 空间中，每一个单点集是一个 $G_{\delta}$ 集。

(b) 存在一个我们熟悉的空间，它的任何一个单点集都是 $G_{\delta}$ 集，但这个空间却不满足第一可数性公理。它是哪个空间呢？

这个术语来源于德文，其中``$G$''就是``Gebiet''，意思是``开集''。``$\delta$''就是``Durchschnitt''，意思是``交''。
\item 证明：若 $X$ 具有可数基 $\{B_{n}\}$，则 $X$ 的每一个基 $\mathcal{C}$ 包含 $X$ 的一个可数基。[提示：对每一对指标 $n,m$，只要可能，便选取 $C_{n,m}\in\mathcal{C}$，使得 $B_{n}\subset C_{n,m}\subset B_{m}$。]
\item 设 $X$ 有可数基，$A$ 是 $X$ 的一个不可数子集。证明 $A$ 中有不可数个点都是 $A$ 的极限点。
\item 证明：每一个紧致度量空间 $X$ 都有可数基。[提示：设 $\mathcal{A}_{n}$ 是由 $1/n$-球构成的 $X$ 的有限覆盖。]
\item (a) 证明：每一个有可数稠密子集的度量空间都有可数基。

(b) 证明：每一个可度量化的 Lindelöf 空间都有可数基。
\item 证明：$\mathbb{R}_{\ell}$ 和 $I_{o}^{2}$ 不可度量化。
\item 对我们所讲的四个可数性公理，空间 $S_{\Omega}$ 满足哪几个？$\overline{S}_{\Omega}$ 又如何？
\item 在一致拓扑下，对我们所讲的四个可数性公理，空间 $\mathbb{R}^{\omega}$ 满足哪几个？
\item 设 $A$ 是 $X$ 的一个闭子空间。证明：如果 $X$ 是 Lindelöf 的，那么 $A$ 也是 Lindelöf 的。举例说明，$X$ 有可数稠密子集，$A$ 却未必有可数稠密子集。
\item 若 $X$ 是可数个有可数稠密子集的空间的积空间，则 $X$ 也有可数稠密子集。
\item 设 $f\colon X\to Y$ 是连续的。证明：如果 $X$ 是 Lindelöf 的，或者有可数稠密子集，那么 $f(X)$ 也满足同样的条件。
\item 设 $f\colon X\to Y$ 是一个连续开映射。证明：如果 $X$ 满足第一或第二可数性公理，那么 $f(X)$ 也满足同样的公理。
\item 证明：若 $X$ 有可数稠密子集，则 $X$ 中的两两无交的开集的族是可数的。
\item 证明：若 $X$ 是 Lindelöf 空间，$Y$ 是紧致的，则 $X\times Y$ 是 Lindelöf 的。
\item 赋予 $\mathbb{R}^{I}$ 一致度量，其中 $I=[0,1]$。设 $\mathcal{C}(I,\mathbb{R})$ 是连续函数空间。证明：$\mathcal{C}(I,\mathbb{R})$ 有可数稠密子集，因此有可数基。[提示：考虑那些连续函数，它们的图形由有限多个线段构成，并且每一个线段的端点是有理数。]
\item (a) 证明：积空间 $\mathbb{R}^{I}$ 有可数稠密子集，其中 $I=[0,1]$。

(b) 证明：如果 $J$ 的基数大于 $\mathcal{P}(\mathbb{Z}_{+})$，那么积空间 $\mathbb{R}^{J}$ 没有可数稠密子集。[提示：若 $D$ 在 $\mathbb{R}^{J}$ 中稠密，定义 $f\colon J\to\mathcal{P}(D)$ 为 $f(\alpha)=D\cap\pi_{\alpha}^{-1}((a,b))$，其中 $(a,b)$ 是 $\mathbb{R}$ 中的一个固定区间。]
\item 赋予 $\mathbb{R}^{\omega}$ 箱拓扑。设 $\mathbb{Q}^{\infty}$ 是由终端为 $0$ 的有理数序列构成的子空间。那么我们所述的四个可数性公理中，这个空间满足哪几个？
\item 设 $G$ 是第一可数的拓扑群。证明：如果 $G$ 有可数稠密子集，或者是 Lindelöf 的，那么 $G$ 有可数基。[提示：设 $\{B_{n}\}$ 是 $e$ 点处的可数基。如果 $D$ 是 $G$ 的可数稠密子集，证明：对于 $d\in D$，集合 $dB_{n}$ 构成了 $G$ 的一个基。如果 $G$ 是 Lindelöf 的，那么对于每一个 $n$，选取一个可数集 $C_{n}$，使得集合 $cB_{n}$ 覆盖 $G$，其中 $c\in C_{n}$。证明：当 $n$ 取遍 $\mathbb{Z}_{+}$ 时，这些集合构成了 $G$ 的一个基。]
\end{enumerate}